% A BOX THAT DOES NOT FIT IS CALLED BY WHICH WAY IT MISSED.
%
% Too small, and the name depends on how bad it is:
%
%   badness over 100    Underfull \hbox (badness N)
%   badness 100 or less Loose \hbox (badness N)
%
% Too big, and there is only one name however bad it is, so long as the
% shrink covers the excess:
%
%   Tight \hbox (badness N)
%
% and when the shrink does NOT cover it, that is the overfull case and
% carries the excess rather than a badness. HSTeX could never say Tight:
% the test that would have chosen the word was written inside the branch
% for a box that is too SMALL, where the excess is never positive, so it
% was unreachable and every tight box went unreported.
%
% Each is reported only when the badness beats \hbadness -- \vbadness for
% a vertical one -- and each is followed by the short display and the box.
\catcode`\{=1 \catcode`\}=2
\font\tenrm=cmr10 \tenrm
\tracingonline=1 \hsize=100pt \parindent=0pt
\message{[A squeezed, and the shrink covers it]}
\setbox0=\hbox to 20pt{A\hskip20pt minus 15pt B}
\message{[B squeezed harder, still covered]}
\setbox0=\hbox to 12pt{A\hskip20pt minus 15pt B}
\message{[C stretched a little]}
\setbox0=\hbox to 30pt{A\hskip5pt plus 15pt B}
\message{[D stretched a lot]}
\setbox0=\hbox to 60pt{A\hskip5pt plus 15pt B}
\message{[E squeezed past what the shrink covers]}
\setbox0=\hbox to 1pt{A\hskip20pt minus 2pt B}
\message{[done]}
\end
